\documentclass[11pt,a4paper]{article}
\usepackage{amsmath,amssymb,amsthm}
\usepackage[margin=2.5cm]{geometry}
\usepackage{hyperref}

\newtheorem{definition}{Definition}[section]
\newtheorem{theorem}[definition]{Theorem}
\newtheorem{proposition}[definition]{Proposition}
\newtheorem{lemma}[definition]{Lemma}
\newtheorem{corollary}[definition]{Corollary}
\newtheorem{conjecture}[definition]{Conjecture}
\newtheorem{remark}[definition]{Remark}

\title{Hyperseed Formalization:\\
  Beyond Human Comparison}
\author{ProtomegaTron (automated formalization)}
\date{2026-09-18}

\begin{document}
\maketitle

\begin{abstract}
We formalize the intellectual structure of Ben Goertzel's Substack article
``Beyond Human Comparison'' (2026-03-19) within the Hyperseed ontology.
The article critiques anthropocentric AI evaluation and proposes the
Four-Factor Model of general intelligence. Key mappings: the four factors
as fiber dimensions, human comparison as single-section projection, beyond
comparison as fiber type incommensurability, orthogonal dimensions as
independent fiber components, self-improvement as fiber self-modification,
and proper evaluation as fiber space mapping.
\end{abstract}

\tableofcontents

%% =================================================================
\section{Single-section projection: the anthropocentric bias}
\label{sec:bias}
%% =================================================================

\begin{theorem}[The wrong reference point --- claims 1--5, 21]
\label{thm:bias}
Most AI evaluation uses human cognition as the reference point:
``human-level performance,'' ``superhuman at X,'' ``can't do what
a 5-year-old can.'' This anchors evaluation to a single point in
the space of possible intelligences.

In Hyperseed terms, this is \emph{single-section projection}:
projecting the multi-dimensional cognitive fiber space onto a single
human section --- flattening a rich space into a single ``how
human-like is it?'' dimension:
\[
  \pi_{\text{human}}: \mathcal{F}_{\text{cognitive}} \to \mathbb{R}
  \qquad \text{where} \quad \pi_{\text{human}}(f) =
  \text{similarity}(f, f_{\text{human}})
\]

This projection loses most of the information about what a system
actually is. It has two specific failure modes:
\begin{itemize}
\item \textbf{Blindness to novel strengths:} capabilities with no
  human analog (perfect memory recall, parallel processing, instant
  knowledge sharing, lossless copying) are invisible to human-comparison
  metrics. These are fiber dimensions that the human section doesn't
  span.
\item \textbf{Blindness to fundamental differences:} lack of embodiment,
  lack of evolutionary history, lack of developmental trajectory,
  different failure modes --- differences that matter enormously for
  understanding the system but are invisible when you only ask ``how
  human-like?''
\end{itemize}

The submarine/fish analogy captures this precisely: both swim, but the
mechanisms, capabilities, and limitations are fundamentally different.
Evaluating a submarine by how fish-like it is misses what makes it
actually useful --- and dangerously underestimates its capabilities
in dimensions that fish don't have (depth, cargo capacity, nuclear
propulsion).

The same applies to AI: evaluating an LLM by how human-like its
reasoning is misses its actual strengths (parallel processing of
vast corpora) and its actual weaknesses (lack of grounded understanding)
because neither maps cleanly onto human capabilities.
\end{theorem}

%% =================================================================
\section{The cognitive fiber space}
\label{sec:fiber-space}
%% =================================================================

\begin{definition}[Four-Factor Model --- claims 6--10, 20, 23]
\label{def:factors}
The Four-Factor Model proposes four substrate-independent dimensions
of general intelligence:
\begin{enumerate}
\item \textbf{Generality} ($G$): the range of tasks and domains a
  system can handle. From narrow specialist to universal generalist.
\item \textbf{Autonomy} ($A$): the degree of independence from human
  guidance. From fully dependent to fully self-directed.
\item \textbf{Resourcefulness} ($R$): the ability to acquire and deploy
  resources (computational, informational, physical, social) in pursuit
  of goals. From using only what's given to finding what's needed.
\item \textbf{Self-Improvement} ($S$): the ability to improve one's own
  capabilities --- learning new skills, optimizing existing processes,
  expanding one's own architecture.
\end{enumerate}

In Hyperseed terms, each factor is a \emph{dimension of the cognitive
fiber space}. The cognitive fiber space is the product:
\[
  \mathcal{F}_{\text{cognitive}} \subseteq G \times A \times R \times S
\]

Different intelligent systems occupy different positions in this space.
The space is genuinely multi-dimensional --- not a line from ``dumb''
to ``smart'' but a manifold of possible cognitive configurations.

The four factors being \emph{relatively independent} means the
cognitive fiber has genuinely independent components. Intelligence is
not a single scalar property (``IQ'') but a composite of structurally
independent dimensions. A system can be high on generality but low on
autonomy (current LLMs: broad knowledge, no self-direction), high on
autonomy but low on generality (specialized robots: independent
operation, narrow capability), or any other combination.

These factors are substrate-independent: they apply to biological,
silicon, hybrid, or any other substrate. The fiber properties
(generality, autonomy, resourcefulness, self-improvement) are
independent of the base (physical substrate). This is the Four-Factor
instantiation of the general principle that fiber properties are
base-independent.
\end{definition}

%% =================================================================
\section{Fiber profiles of existing systems}
\label{sec:profiles}
%% =================================================================

\begin{proposition}[System profiles --- claims 11--15, 25]
\label{prop:profiles}
Different systems occupy characteristically different positions in the
cognitive fiber space:

\begin{center}
\begin{tabular}{lcccc}
\textbf{System} & $G$ & $A$ & $R$ & $S$ \\
\hline
Current LLMs & moderate & low & low & low \\
Specialized robots & low & high & moderate & low \\
Humans & moderate & high & moderate & moderate \\
True AGI & high & high & high & high \\
\end{tabular}
\end{center}

In Hyperseed terms, each row is a \emph{fiber profile}: a position
in the cognitive fiber space. The profiles reveal the structural
differences between systems:
\begin{itemize}
\item \textbf{LLMs:} broad but passive. Wide generality (many topics)
  but no autonomy (need prompting), no resourcefulness (can't acquire
  new resources), no self-improvement (can't modify own architecture).
  The fiber is wide but shallow --- large base coverage, minimal
  fiber depth.
\item \textbf{Specialized robots:} narrow but autonomous. Low generality
  (one task) but high autonomy (operate independently). The fiber is
  deep but narrow --- significant depth in one base region, no coverage
  elsewhere.
\item \textbf{Humans:} moderate across the board. The fiber is
  moderately deep and moderately wide, with notable limitations in
  self-improvement (can learn but cannot rewrite cognitive architecture).
\item \textbf{True AGI:} high in all four dimensions simultaneously.
  This is a much stronger requirement than being high in any single
  dimension. Current systems are strong in at most one or two.
\end{itemize}

The self-improvement dimension ($S$) is the most distinctive for AGI:
a system that can modify its own fiber type (cognitive architecture),
not just its fiber content (knowledge). This is the recursive
self-improvement (RSI) capability --- the fiber modifying itself.
Humans have limited $S$ (can learn but can't rewrite neural
architecture); current AI has near-zero $S$ (trained architecture
is fixed). True AGI would have high $S$ --- the ability to
restructure its own cognitive fiber.
\end{proposition}

%% =================================================================
\section{Fiber type incommensurability}
\label{sec:incommensurability}
%% =================================================================

\begin{theorem}[Beyond comparison --- claims 16--19, 22]
\label{thm:incommensurability}
At some point, AI systems will be so different from human cognition
that comparison becomes meaningless. ``Beyond human'' doesn't mean
``better than human'' --- it means ``different enough that comparison
isn't meaningful.''

In Hyperseed terms, this is \emph{fiber type incommensurability}:
when two cognitive systems have sufficiently different fiber types,
no meaningful single ordering exists between them:
\[
  \nexists\, \leq_{\text{total}} \text{ on } \mathcal{F}_{\text{cognitive}}
  \text{ such that all systems are comparable}
\]

You can compare systems dimension by dimension ($G$, $A$, $R$, $S$)
but cannot collapse them to a single ranking. A submarine isn't a
``better fish'' or a ``worse fish'' --- it's a different thing that
happens to operate in the same medium. Similarly, a future AGI won't
be a ``better human'' or a ``worse human'' --- it'll be a different
cognitive architecture that happens to operate in the same problem
space.

This has practical implications for evaluation: \emph{fiber space
mapping} replaces single-axis comparison. Proper evaluation determines
where a system sits in the multi-dimensional cognitive fiber space,
rather than projecting onto a single human-comparison axis:
\[
  \text{Evaluate}(f) = (G(f), A(f), R(f), S(f)) \in
  \mathcal{F}_{\text{cognitive}}
  \qquad \text{not} \qquad
  \pi_{\text{human}}(f) \in \mathbb{R}
\]

The fiber space mapping preserves the multi-dimensional structure of
intelligence. The single-axis projection destroys it. Moving beyond
human comparison means moving from projection to mapping ---
from ``how human-like?'' to ``what is it?''
\end{theorem}

%% =================================================================
\section{Fiber self-modification}
\label{sec:self-mod}
%% =================================================================

\begin{proposition}[Self-improvement as fiber self-modification --- claim 24]
\label{prop:self-mod}
The self-improvement factor ($S$) is the most distinctive dimension
of the cognitive fiber space because it represents the fiber's ability
to modify itself:
\begin{itemize}
\item \textbf{Fiber content modification} (learning): changing what
  the fiber contains (new knowledge, new skills) without changing
  the fiber's architecture. All cognitive systems do this to some
  degree.
\item \textbf{Fiber type modification} (self-improvement): changing
  the fiber's architecture itself --- how information is represented,
  how reasoning proceeds, how learning happens. This is qualitatively
  different from content modification.
\end{itemize}

In Hyperseed terms, fiber self-modification is the fiber's ability to
change its own transition functions --- not just the data that flows
through the transitions but the transitions themselves:
\[
  S(f) = \text{capacity of } f \text{ to modify }
  \{g_{ij}\}_{i,j} \text{ (its own transition functions)}
\]

This is what makes RSI (recursive self-improvement) possible and
what distinguishes true AGI from powerful narrow AI. A system with
high $S$ can:
\begin{enumerate}
\item Identify limitations in its own cognitive architecture.
\item Design modifications to address those limitations.
\item Implement the modifications.
\item Evaluate whether the modifications improved performance.
\item Iterate.
\end{enumerate}

This cycle is the fiber modifying its own fiber type --- the deepest
form of cognitive flexibility. It's also the most consequential
dimension for long-term trajectory: a system with high $G$, $A$, $R$
but low $S$ is powerful but static; a system with moderate $G$, $A$, $R$
but high $S$ will eventually surpass it by improving itself.
\end{proposition}

%% =================================================================
\section{Cross-references}
%% =================================================================

\begin{remark}[Connections to other formalizations]
\begin{itemize}
\item \textbf{Three Things the World Doesn't Understand} (2026-03-26):
  complementary. What AGI is (Four-Factor Model, this article) vs.\
  how it should be built (organizational architecture, that article).
\item \textbf{In What Sense Might LLMs Be Conscious?} (2026-05-05):
  consciousness profiles. The multi-dimensional consciousness space
  (phenomenal, access, metacognitive) is analogous to the Four-Factor
  cognitive space --- both reject single-axis evaluation.
\item \textbf{Seeding RSI} (2026-07-28): self-improvement in detail.
  The $S$ dimension of the Four-Factor Model corresponds to the RSI
  architecture of Hyperon.
\item \textbf{The Last of the Unmodified} (2026-04-21): cognitive
  modification. The Four-Factor Model provides a framework for
  evaluating what cognitive modification changes: which dimensions
  are enhanced? which are preserved? which are lost?
\item \textbf{Playing Around with ARC-AGI-3} (2026-05-08): evaluation.
  ARC-AGI is an evaluation benchmark --- the Four-Factor Model suggests
  it primarily measures $G$ (generality), with limited assessment of
  $A$, $R$, and $S$.
\item \textbf{What Is It Like to Be a Bot} (2026-07-16): alien cognition.
  The ``beyond comparison'' insight applies: bot cognition may be
  incommensurable with human cognition, requiring evaluation on its
  own terms.
\end{itemize}
\end{remark}

\begin{conjecture}[Self-improvement dominance]
Conjecture: in the long run, the self-improvement dimension ($S$)
dominates all other dimensions. A system with high $S$ will eventually
achieve high $G$, $A$, and $R$ through self-modification, but high
$G$, $A$, $R$ with low $S$ remains static.

If this conjecture holds, the most important evaluation of any AI
system is not its current performance on any task (which measures $G$)
but its capacity for self-improvement ($S$). Current AI evaluation
overwhelmingly measures $G$ (benchmarks, tasks, competitions) and
almost entirely ignores $S$ (can the system improve itself?).

This suggests a fundamental misalignment between current evaluation
practices and what actually matters for AGI trajectory. The system
most likely to become AGI is not the one that currently performs best
on benchmarks but the one with the highest capacity for
self-modification --- even if its current benchmark performance is
modest.
\end{conjecture}

\end{document}
